\documentclass[12pt,a4paper]{article}

% =========================
% Pacotes essenciais
% =========================
\usepackage[a4paper,margin=2.5cm]{geometry}
\usepackage[T1]{fontenc}
\usepackage[utf8]{inputenc}
\usepackage[brazil]{babel}
\usepackage{lmodern}
\usepackage{setspace}
\usepackage{microtype}
\usepackage{hyperref}
\usepackage{enumitem}
\usepackage{fancyhdr}

\hypersetup{
  colorlinks=true,
  linkcolor=black,
  urlcolor=blue
}

\setstretch{1.15}
\setlength{\parindent}{0pt}
\setlength{\parskip}{6pt}

\pagestyle{fancy}
\fancyhf{}
\lhead{Manual do Aluno — CV721A}
\rhead{\thepage}
\cfoot{}

% =========================
% Metadados
% =========================
\newcommand{\Disciplina}{Fundações (CV721A)}
\newcommand{\Curso}{Engenharia Civil — UNICAMP}
\newcommand{\VersaoManual}{v1.0}
\newcommand{\DataManual}{\today}

\title{\textbf{Manual do Aluno}\\
Sistema de Presença por QR Code\\
\Disciplina}
\author{}
\date{\DataManual\ (\VersaoManual)}

\begin{document}
\maketitle

\section{Apresentação}

Este manual explica como utilizar o \textbf{Sistema de Presença por QR Code}, adotado na disciplina \Disciplina{}, do curso de \Curso{}.

O objetivo do sistema é tornar o registro de presença \textbf{rápido, simples e organizado}, evitando listas em papel e reduzindo erros de registro.

---

\section{O que é o QR Code de Presença}

Em cada aula presencial, o professor exibirá um \textbf{QR Code} no início da aula.

Esse QR Code:
\begin{itemize}[noitemsep]
  \item é válido apenas por alguns minutos;
  \item deve ser acessado \textbf{durante a aula};
  \item substitui a lista de presença tradicional.
\end{itemize}

---

\section{Como Registrar sua Presença}

Siga os passos abaixo:

\begin{enumerate}[label=\arabic*)]
  \item Abra a câmera do seu celular ou um aplicativo leitor de QR Code.
  \item Aponte a câmera para o QR Code exibido pelo professor.
  \item Clique no link que aparecer na tela.
  \item Preencha corretamente:
  \begin{itemize}[noitemsep]
    \item seu \textbf{RA};
    \item seu \textbf{Nome completo}.
  \end{itemize}
  \item Clique em \textbf{Confirmar presença}.
\end{enumerate}

Após a confirmação, sua presença estará registrada no sistema.

---

\section{Regras Importantes}

\begin{itemize}[noitemsep]
  \item O QR Code tem \textbf{tempo limitado}. Registre sua presença assim que ele for exibido.
  \item O registro deve ser feito \textbf{pelo próprio aluno}.
  \item Cada RA pode registrar presença \textbf{apenas uma vez por aula}.
  \item Registros fora do horário ou tentativas duplicadas não serão aceitos.
\end{itemize}

---

\section{Problemas Comuns}

\subsection*{Não consegui registrar a presença}
Verifique:
\begin{itemize}[noitemsep]
  \item se o QR Code ainda está válido;
  \item se seu celular está conectado à internet;
  \item se o RA e o nome foram digitados corretamente.
\end{itemize}

Caso o problema persista, informe o professor \textbf{ainda durante a aula}.

---

\section{Uso dos Dados}

Os dados coletados (RA, nome, data e horário) são utilizados exclusivamente para:
\begin{itemize}[noitemsep]
  \item controle de frequência da disciplina;
  \item consolidação de informações acadêmicas.
\end{itemize}

Os dados seguem as normas institucionais da UNICAMP e não são utilizados para outros fins.

---

\section{Considerações Finais}

O Sistema de Presença por QR Code foi desenvolvido para facilitar o acompanhamento da disciplina e tornar o processo mais transparente e organizado.

A colaboração dos alunos é essencial para o bom funcionamento do sistema.

\end{document}

